\chapter{Perangkat Keras}
\label{ch:03-perangkat-keras}

\section{Gambaran Blok Perangkat Keras}

Node perangkat terdiri atas satu papan pengendali dan empat modul: sensor
ultrasonik sebagai pengukur jarak, motor servo sebagai penggerak gate pakan, modul
RTC sebagai penjaga waktu, dan LCD karakter sebagai penampil lokal.
\figref{fig:hw-block-diagram} menunjukkan hubungan antar blok beserta jenis
antarmukanya. Semua modul berbagi ground yang sama; yang berbeda adalah jalur
antarmukanya: sensor memakai dua GPIO diskret, servo memakai satu kanal PWM,
sedangkan RTC dan LCD berbagi satu bus I\textsuperscript{2}C.

\figuretikz[\linewidth]{hw-block-diagram}{Blok perangkat keras node: papan \BoardName{} dengan
sensor ultrasonik, motor servo, modul RTC, dan LCD karakter beserta jenis
antarmukanya.}{fig:hw-block-diagram}

Setiap modul diperiksa keberadaannya saat \code{setup()} dan setiap fungsi
perangkat keras mengembalikan nilai yang dapat diperiksa pemanggilnya, sehingga
modul yang tidak terdeteksi hanya menghilangkan kemampuan terkait, bukan
menghentikan program (\chapref{ch:04-firmware}).

\section{Daftar Komponen}

\tabref{tab:03-bom} memuat komponen prototipe beserta tegangan kerja dan
antarmukanya. Kolom catatan memuat hal yang paling mudah salah saat perakitan.

\begin{table}[H]
\centering
\caption{Daftar komponen node perangkat.}
\label{tab:03-bom}
\small
\begin{tabular}{@{}p{3.2cm}p{1.8cm}p{3.1cm}p{6.3cm}@{}}
\toprule
\textbf{Komponen} & \textbf{Tegangan} & \textbf{Antarmuka} & \textbf{Catatan} \\
\midrule
\BoardName & 5 V (USB) & USB, WiFi 2.4 GHz & Pengendali utama; regulator papan
menyuplai 3.3 V untuk modul I\textsuperscript{2}C. \\
Sensor ultrasonik HC-SR04 & 5 V & 2 GPIO diskret & \code{TRIG} menerima 3.3 V;
\code{ECHO} mengeluarkan 5 V sehingga wajib dibagi tegangan. \\
Motor servo SG90 & 5 V & PWM \ServoFreq{} Hz & Jangan diberi 3.3 V; arus stall
besar saat mulai bergerak. \\
Modul RTC DS3231 & 3.3 V & I\textsuperscript{2}C \AddrRtc & Butuh sel koin CR2032
agar waktu bertahan saat daya putus. \\
LCD1602 + backpack PCF8574 & 3.3 V & I\textsuperscript{2}C \AddrLcd{} atau
\AddrLcdAlt & Alamat bergantung jumper backpack; dipindai otomatis. \\
Resistor pembagi & -- & 1 k$\Omega$ dan 2 k$\Omega$ & Menurunkan \code{ECHO} 5 V
menjadi sekitar 3.3 V. \\
Kapasitor elektrolit & 5 V & 470--1000 µF & Melintang VCC dan GND servo sebagai
penahan lonjakan arus. \\
Sel koin CR2032 & 3 V & Soket DS3231 & Tanpa sel, jam kembali tidak valid setiap
kali daya putus. \\
\bottomrule
\end{tabular}
\end{table}

\src{firmware/esp32_feeder/README.md}

\section{Skema Rangkaian}

\figref{fig:hw-schematic} menggambarkan sambungan listrik yang harus dibuat. Tiga
hal perlu ditekankan karena ketiganya adalah sumber kegagalan yang paling sering
muncul pada prototipe: pembagi tegangan pada \code{ECHO}, sumber daya servo, dan
ground bersama.

\figuretikz[\linewidth]{hw-schematic}{Skema rangkaian: pembagi tegangan pada jalur \code{ECHO},
bus I\textsuperscript{2}C untuk DS3231 dan LCD, serta jalur daya servo.}{fig:hw-schematic}

\subsection{Pembagi Tegangan pada ECHO}

HC-SR04 bekerja pada 5 V dan mengeluarkan pulsa \code{ECHO} setinggi 5 V, sedangkan
GPIO ESP32 tidak toleran terhadap 5 V. Karena itu jalur \code{ECHO} tidak boleh
disambung langsung ke \PinEcho. Rangkaian yang dipakai adalah pembagi sederhana:
\code{ECHO} melewati resistor 1 k$\Omega$ menuju titik sambung, dari titik itu
resistor 2 k$\Omega$ ke ground, dan titik sambung dihubungkan ke GPIO \PinEcho.
Dengan rasio itu tegangan 5 V turun menjadi sekitar 3.3 V. Gejala khas bila pembagi
ini terlewat adalah jarak yang selalu tidak valid atau pin yang rusak setelah
beberapa lama pemakaian.

\begin{warnbox}
\code{ECHO} pada HC-SR04 adalah keluaran 5 V, sedangkan GPIO \PinEcho{} tidak
\emph{5 V tolerant}. Pembagi 1 k$\Omega$ / 2 k$\Omega$ bukan pilihan opsional,
melainkan bagian dari rangkaian.
\end{warnbox}

\subsection{Daya Servo}

SG90 menarik arus kecil saat diam, tetapi arus stall-nya sekitar 700 mA pada 5 V dan
lonjakan itu terjadi tepat saat gate mulai membuka atau menutup. Bila servo diberi
daya dari jalur 5 V papan yang juga memberi daya ESP32 lewat USB, lonjakan tersebut
menurunkan tegangan suplai sesaat dan papan dapat melakukan brownout reset. Gejalanya
khas: papan restart setiap kali servo bergerak dan log serial mengulang urutan
inisialisasi. Mitigasinya adalah suplai 5 V terpisah dengan ground disatukan, serta
kapasitor elektrolit 470--1000 µF melintang VCC dan GND servo. Menghubungkan VCC
servo ke pin 3.3 V bukan mitigasi, melainkan kesalahan yang lebih parah.

\subsection{Bus I\textsuperscript{2}C}

DS3231 dan LCD berbagi SDA pada \PinSda{} dan SCL pada \PinScl. Bus berjalan pada
\IicClock{} karena firmware tidak memanggil \code{Wire.setClock()}, sehingga
kecepatan mengikuti nilai bawaan pustaka. Modul DS3231 dan backpack PCF8574 umumnya
sudah menyertakan resistor pull-up pada papannya; bila keduanya terpasang, kedua
pasangan pull-up berada paralel dan bus tetap bekerja. Pull-up tambahan hanya
diperlukan bila salah satu modul berupa papan telanjang.

\section{Pembagian Pin}

\figref{fig:hw-pinout} memetakan pin yang dipakai pada papan 30 pin, dan
\tabref{tab:03-pin} merinci makro firmware, GPIO, fungsi, arah, level, serta
catatannya. Tabel ini adalah padanan langsung dari \file{config.h}, sehingga
perubahan pin cukup dilakukan di satu tempat itu.

\figuretikz[\linewidth]{hw-pinout}{Peta pin papan \BoardName{} yang dipakai node perangkat.}{fig:hw-pinout}

\begin{table}[H]
\centering
\caption{Pembagian pin node perangkat.}
\label{tab:03-pin}
\small
\begin{tabular}{@{}p{2.0cm}p{1.1cm}p{2.6cm}p{1.2cm}p{1.9cm}p{4.6cm}@{}}
\toprule
\textbf{Makro} & \textbf{GPIO} & \textbf{Fungsi} & \textbf{Arah} & \textbf{Level} &
\textbf{Catatan} \\
\midrule
\code{PIN_SERVO} & \PinServo & PWM SG90 & keluar & 3.3 V & Pulsa
\ServoMinUs--\ServoMaxUs{} µs. \\
\code{PIN_TRIG} & \PinTrig & Trigger HC-SR04 & keluar & 3.3 V & Diberi nilai rendah
saat inisialisasi. \\
\code{PIN_ECHO} & \PinEcho & Echo HC-SR04 & masuk & 5 V dibagi & Wajib melalui
pembagi 1 k$\Omega$ / 2 k$\Omega$. \\
\code{PIN_SDA} & \PinSda & I\textsuperscript{2}C data & dua arah & 3.3 V &
Open-drain; DS3231 dan LCD. \\
\code{PIN_SCL} & \PinScl & I\textsuperscript{2}C clock & keluar & 3.3 V & Jam bus
\IicClock. \\
\bottomrule
\end{tabular}
\end{table}

\src{firmware/esp32_feeder/config.h}

\tabref{tab:03-pin-terlarang} memuat pin yang tidak boleh dipakai, baik karena sudah
dipakai fungsi internal papan maupun karena tidak tersedia pada modul 30 pin.
Daftar ini bukan catatan teoretis: memakai GPIO 6 sampai 11, misalnya, akan
menghalangi papan melakukan boot karena pin tersebut terhubung ke flash SPI.

\begin{table}[H]
\centering
\caption{Pin yang tidak boleh dipakai pada papan 30 pin.}
\label{tab:03-pin-terlarang}
\small
\begin{tabular}{@{}p{1.8cm}p{2.8cm}p{9.9cm}@{}}
\toprule
\textbf{GPIO} & \textbf{Status} & \textbf{Alasan} \\
\midrule
6--11 & terlarang & Terhubung ke flash SPI internal yang dipakai papan membaca
program. \\
1, 3 & terlarang & UART0 yang dipakai jembatan USB-serial CH340 untuk unggah program
dan log serial. \\
0, 2, 12, 15 & terlarang & Pin boot strapping; GPIO 2 sekaligus terhubung ke LED
bawaan papan. \\
34--39 & terlarang & Masukan saja, tanpa pull-up internal dan tanpa keluaran. \\
20, 24, 28--31 & tidak tersedia & Tidak dibawa keluar pada modul 30 pin. \\
2, 4, 12, 14, 15, 16, 17, 19, 23, 25, 26, 27, 32, 33, 34, 35, 36, 39 & tidak
dipakai & Sisa pin modul yang bebas untuk penambahan, dengan memperhatikan baris di
atas. \\
\bottomrule
\end{tabular}
\end{table}

\section{Timing Servo}

Gate digerakkan sebagai urutan buka, tutup, jeda, dan ulang, bukan satu gerakan
panjang. \tabref{tab:03-timing} merangkum parameter waktunya; semuanya diambil dari
\file{config.h}.

\begin{table}[H]
\centering
\caption{Parameter timing gate pakan dan pensinyalan servo.}
\label{tab:03-timing}
\small
\begin{tabular}{@{}p{4.6cm}p{2.4cm}p{7.5cm}@{}}
\toprule
\textbf{Parameter} & \textbf{Nilai} & \textbf{Arti} \\
\midrule
\code{PORTION_OPEN_MS} & \PortionOpenMs{} ms & Lama gate terbuka untuk satu porsi. \\
\code{SERVO_SETTLE_MS} & \SettleMs{} ms & Jeda setelah gate ditutup agar mekanik
berhenti. \\
\code{PORTION_GAP_MS} & \GapMs{} ms & Jeda antar porsi dalam satu perintah. \\
Siklus per porsi & \PortionCycleMs{} ms & Waktu buka ditambah settle; satuan durasi
total. \\
\code{SERVO_MIN_US} & \ServoMinUs{} µs & Lebar pulsa untuk sudut minimum. \\
\code{SERVO_MAX_US} & \ServoMaxUs{} µs & Lebar pulsa untuk sudut maksimum. \\
Frekuensi PWM & \ServoFreq{} Hz & Periode 20 ms, sesuai servo analog. \\
Sudut tertutup & \ServoClosedDeg\textdegree & Posisi gate menutup aliran pakan. \\
Sudut terbuka & \ServoOpenDeg\textdegree & Posisi gate membuka aliran pakan. \\
\bottomrule
\end{tabular}
\end{table}

\src{firmware/esp32_feeder/config.h}

Durasi satu perintah pakan adalah $N \times \PortionCycleMs + (N-1) \times \GapMs$
milidetik untuk $N$ porsi: satu porsi memerlukan \PortionCycleMs{} ms, dua porsi
menambahkan satu jeda \GapMs{} ms, dan sepuluh porsi menambahkan sembilan jeda.
Angka itu batas bawah, karena setelah gate terakhir menutup firmware masih mengukur
ulang jarak, mempublikasikan telemetri, lalu mengirim ack. Karena itu ack selalu
mencerminkan keadaan hopper sesudah pemberian pakan, bukan sebelum.

\section{Kalibrasi Hopper}

Sensor mengukur jarak dari mulut hopper ke permukaan pakan, sehingga nilainya
berbanding terbalik dengan isi hopper. Konversi memakai dua titik acuan, yaitu
\HopperFull{} cm saat penuh dan \HopperEmpty{} cm saat kosong, dengan rentang
efektif \HopperSpan{} cm. Rumus yang dipakai firmware maupun platform adalah

\begin{center}
\PctFormula
\end{center}

dengan $d$ jarak terukur dalam sentimeter; hasilnya dibatasi 0 sampai 100. Untuk
jarak 10 cm, selisih terhadap jarak kosong adalah 7 cm dari rentang \HopperSpan{}
cm, sehingga hasilnya \PctAtTenCm{} persen. Angka itu penting sebagai pembanding:
dengan kalibrasi lain, jarak yang sama menghasilkan persentase yang berbeda jauh.
Karena itu nilai \HopperFull{} dan \HopperEmpty{} di firmware, di skema basis data,
dan di baris \code{devices} dijaga tetap sama; platform memakai nilai di basis data,
firmware memakai nilai di \file{config.h}, dan keduanya harus diubah bersama.

Pembacaan di atas \MaxDistance{} cm dianggap tidak valid karena di luar rentang itu
pantulan ultrasonik tidak dapat dipercaya. Pada keadaan itu firmware mengirim
telemetri berstatus \code{error} tanpa field \code{distance_cm}, dan platform
menyimpan baris telemetri tanpa persentase. Pembacaan tidak valid tidak boleh
diperlakukan sebagai nol persen, karena nol persen berarti hopper kosong dan memicu
notifikasi yang berbeda.
